\section{Governance Failures in Multi-Tenant Platforms}
\label{sec:problem}

Multi-tenant platforms usually start with a clean story. A tenant owns data. Users belong to tenants. Roles grant capabilities. Services enforce access. The model is fine. The implementation is where it frays.

\subsection{Scope ambiguity}

A service receives a tenant identifier from a caller and treats it as truth. A request header says \code{tenant=acme}; a query parameter says \code{schema=demo_acme}; a job payload carries an old prefix. Each value is plausible. None is authoritative unless it was derived from authenticated context by the system that owns tenancy.

That is the confused-deputy shape in platform clothing \cite{confused_deputy}. A downstream service acts with more authority than the caller should have, or under a tenant the actor did not intend, because the service trusted a scope hint it should have treated as untrusted input.

\subsection{Governance scatter}

The second failure is duplication. Every service that joins the platform implements its own version of identity resolution, module gating, role checks, impersonation rules, and audit logging. One service validates JWTs against Cognito. Another trusts an internal caller. Another reads a tenant table directly. Another copies old middleware and never gets the patch.

At 12 services, a single governance rule change becomes 12 separate edits, 12 review paths, and 12 chances to drift. The code volume is not the main problem. The problem is that the platform no longer has one answer.

\subsection{Surface drift}

The third failure is what the platform advertises. API documentation and generated clients often expose every endpoint the product knows how to serve, not every endpoint a given tenant may use. Clients discover disabled modules, call their routes, and receive errors that look like product bugs. Operators see support tickets. Security teams see unnecessary attack surface.

The visible API surface should be a consequence of effective tenant capability. If a tenant has not licensed a module, the API spec should not invite them to call it.

\subsection{Design goals}

The two-plane boundary is designed around five goals.

\begin{description}[leftmargin=2.2em, style=nextline]
  \item[G1 Explicit tenancy.] Actor identity and effective tenant context are distinct fields. Support impersonation changes the effective context, never the actor.
  \item[G2 Scope injection.] Scope flows downstream from the control plane. It does not flow upstream from callers.
  \item[G3 Minimized surface.] Tenants see only the routes and modules they are allowed to use.
  \item[G4 Auditable impersonation.] Operators can view as a tenant without shared credentials or unattributed actions.
  \item[G5 Governance as data.] Routine tenant operations become reviewed configuration changes. New module types still require code.
\end{description}

These goals are modest on purpose. The point is not to build a universal policy engine. The point is to remove ambiguity from the high-frequency paths that make SaaS platforms painful to operate.
